% ============================================================
% SECTION II: RELATED WORK
% ============================================================
\section{Related Work}

We walk through the deepfake-detection literature roughly in the order in which the field developed, starting from the paper that gave us the standard benchmark and the Xception baseline, and ending with the most recent transformer-based work. The arc is short and consistent: binary detection on a single dataset $\rightarrow$ generalization problems $\rightarrow$ search for forensic signals (frequency, blending) that survive across datasets $\rightarrow$ the harder question of \textit{which} method produced a fake. DSAN sits at the end of this arc.

\subsection{The base paper and the binary-detection era}
The story begins with R\"ossler~\textit{et~al.}~\cite{rossler2019faceforensics}, whose FaceForensics++ (FF++) dataset became the de-facto benchmark for the field. They paired four manipulation methods - Deepfakes, Face2Face, FaceSwap, and NeuralTextures - with 1,000 source videos and showed that an XceptionNet~\cite{chollet2017xception} fine-tuned on raw frames already exceeds 95\% binary accuracy. This is the \textbf{base paper} our work builds on: we keep XceptionNet as the binary spatial gatekeeper, but treat the four FF++ methods as four \textit{attribution classes} rather than a single ``fake'' bucket. The systematic review by Rana~\textit{et~al.}~\cite{rana2022slr}, covering 112 papers, confirms what early adopters already felt: CNN-based detectors dominate the field, and generalization remains the open problem.

\subsection{Looking beyond pixels: biological and frequency cues}
With pixel-level detectors saturating on FF++, researchers started looking for signals that might be more universal. Jung~\textit{et~al.}~\cite{jung2020deepvision} took the biological route, tracking eye-blinking patterns to detect deepfakes and reaching 87.5\% accuracy. The idea is appealing - early generators really did produce unnatural blink behaviour - but newer generators have largely closed that gap, which is one reason we ultimately did not include a blink module in DSAN.

A more durable line of work came from the frequency domain. Frank~\textit{et~al.}~\cite{frank2020leveraging} showed that GAN upsampling leaves periodic spectral artifacts in the 2D Fourier spectrum, visible as a regular grid of peaks. This finding ties into the much older steganalysis work of Fridrich and Kodovsk\'y~\cite{fridrich2012srm}, whose Spatial Rich Model (SRM) noise residuals are still the standard tool for surfacing low-level forensic traces. Hao~\textit{et~al.}~\cite{hao2021twostreamSRM} were the first to combine the two ideas explicitly: a two-stream network that fuses RGB features with learnable SRM residuals for binary detection. Our frequency stream is a direct intellectual descendant of theirs - we keep the SRM input but add 2D FFT magnitude/phase/power channels and, crucially, push it through to attribution rather than stopping at binary detection.

\subsection{The generalization problem}
By 2020-2022 it became obvious that the real challenge was not in-domain accuracy but cross-dataset robustness. Two complementary ideas changed the landscape here. Li~\textit{et~al.}~\cite{li2020facexray} introduced \textbf{Face X-ray}, observing that almost every face-swap or reenactment method leaves a blending boundary, and that learning to localize this boundary gives a manipulation-agnostic signal. Shiohara and Yamasaki~\cite{shiohara2022sbi} pushed this further with \textbf{Self-Blended Images} (SBI): pseudo-fakes generated by self-blending real images through soft masks, which avoids overfitting to any specific generator. We use both ideas in DSAN - SBI as augmentation, and a lightweight Face-X-ray-style mask decoder as an auxiliary loss. Tran~\textit{et~al.}~\cite{tran2023metadeepfake} attacked the generalization gap from a different angle, using meta-learning with pair-attention and alignment losses across multiple source domains; their gains on unseen datasets motivate domain-aware training tricks of the sort we adopt.

\subsection{The recent shift: transformers, fairness, and the meta-question}
Three recent papers reshape how the community now thinks about deepfake detection. Xu~\textit{et~al.}~\cite{xu2024fairness} performed a large-scale fairness audit using massively annotated databases (47 attributes spanning age, gender, ethnicity) and found measurable bias: detection accuracy varies by demographic group. Their finding is sobering for any system claiming forensic readiness, and is one reason our discussion section explicitly flags identity-safe evaluation as non-negotiable. Khan and Dang-Nguyen~\cite{khan2024generalization} ran a controlled study across architectures (CNN vs.\ Transformer), datasets, and pre-training paradigms (supervised, DINO, CLIP). Their headline result: transformer backbones with self-supervised pre-training generalize markedly better than CNNs trained from scratch, and FF++ plus DFDC offer the broadest training diversity - which justifies our choice to train on FF++ and probe on Celeb-DF and DFDC. Khormali and Yuan~\cite{khormali2024selfsupgraph} take the transformer trend one step further with a self-supervised graph-transformer detector that explicitly models relationships between facial regions; they report strong cross-dataset numbers and good robustness to compression and noise. Edwards~\textit{et~al.}~\cite{edwards2024review} provide the most up-to-date review of the area; their two consistent findings - DL-based methods dominate, but generalization is the open problem - frame the position our paper takes.

\subsection{Loss-side and explainability tools we adopt}
DSAN inherits two algorithmic ideas. Khosla~\textit{et~al.}~\cite{khosla2020supervised} introduced \textbf{Supervised Contrastive Learning}, which we adapt with a higher temperature ($\tau{=}0.15$) appropriate to our smaller effective batch size. For explainability we use Grad-CAM++~\cite{chattopadhay2018gradcampp}, but produce \textit{two} heatmaps - one per stream - rather than the usual single spatial map. To our knowledge this dual-domain visualization is novel for attribution.

\subsection{Summary and where DSAN fits}
Table~\ref{tab:related_summary} consolidates the storyline above with each work's core contribution and the limitation that motivates our design. Three patterns emerge:
(i) almost all prior work targets binary detection - none of the surveyed papers reports per-method attribution F1 on FF++ with identity-safe splits;
(ii) frequency-domain forensics is repeatedly shown to help, but is rarely combined with explainability;
(iii) cross-dataset evaluation, when reported, consistently shows a 15-25 point AUC drop, which we observe and report honestly in Section~\ref{sec:cross_dataset}.

DSAN addresses these gaps by (a) reframing the four FF++ methods as an attribution problem with identity-safe splits, (b) fusing RGB and frequency through a per-dimension gate that does not collapse the way two-token attention does, and (c) producing dual Grad-CAM++ maps for each prediction.

\begin{table*}[!t]
\centering
\caption{Storyline summary of related work. Each row is a conceptual milestone in the field's evolution from FF++/Xception (the base paper) to recent transformer / fairness work.}
\label{tab:related_summary}
\renewcommand{\arraystretch}{1.2}
\footnotesize
\begin{tabular}{@{}p{0.6cm} p{2.6cm} p{3.7cm} p{2.6cm} p{2.9cm} p{3.4cm}@{}}
\toprule
\textbf{Year} & \textbf{Work} & \textbf{Approach / Backbone} & \textbf{Dataset(s)} & \textbf{Key finding} & \textbf{Limitation we address} \\
\midrule
2019 & R\"ossler \textit{et~al.}~\cite{rossler2019faceforensics} \textit{(base paper)} & XceptionNet on FF++ frames & FF++ (4 methods) & 95\%+ binary accuracy on c23 & Binary only; no method attribution \\
2020 & Jung \textit{et~al.}~\cite{jung2020deepvision} & Eye-blink heuristic + ML & Custom DF videos & Blink anomalies useful, 87.5\% acc. & Newer generators close the blink gap \\
2020 & Frank \textit{et~al.}~\cite{frank2020leveraging} & 2D FFT spectral analysis & FF++, GAN sets & GAN upsampling leaves grid peaks & Frequency-only; not used for attribution \\
2020 & Li \textit{et~al.}~\cite{li2020facexray} & Face X-ray blending boundary & FF++, DFDC, Celeb-DF & Boundary is a universal manipulation cue & Trained for binary only \\
2021 & Hao \textit{et~al.}~\cite{hao2021twostreamSRM} & RGB + learnable SRM two-stream & FF++ & Frequency complements RGB & Binary only; no FFT, no attribution \\
2022 & Shiohara \& Yamasaki~\cite{shiohara2022sbi} & Self-Blended Images (SBI) & FF++ $\to$ Celeb-DF & Generator-free pseudo-fakes & Augmentation only, no attribution \\
2022 & Rana \textit{et~al.}~\cite{rana2022slr} & Systematic review (112 papers) & Multiple & Generalization is the open problem & Survey, no new method \\
2023 & Tran \textit{et~al.}~\cite{tran2023metadeepfake} & Meta-learning (MDD) + alignment & Multi-domain splits & Better unseen-domain generalization & Binary only; expensive meta-training \\
2024 & Xu \textit{et~al.}~\cite{xu2024fairness} & Fairness audit (47 attributes) & Annotated FF++/DFDC subsets & Detector bias varies by demographic & Forces identity-safe / fair eval \\
2024 & Khan \& Dang-Nguyen~\cite{khan2024generalization} & CNN vs.\ Transformer + DINO/CLIP & FF++, DFDC, Celeb-DF, FakeAVCeleb & Transformers + SSL generalize best & CNN backbones still common in attribution \\
2024 & Khormali \& Yuan~\cite{khormali2024selfsupgraph} & Self-supervised graph transformer & Cross-dataset benchmarks & Strong robustness to noise/compression & Binary only; high compute \\
\midrule
\textbf{Ours} & \textbf{DSAN (this work)} & \textbf{Dual-stream gated fusion + SupCon + mask + dual Grad-CAM++} & \textbf{FF++ (4-class attribution)} & \textbf{Per-method attribution with explainability} & \textbf{Addresses (i)--(iii) above} \\
\bottomrule
\end{tabular}
\end{table*}
